\documentclass[11pt]{article}
\usepackage{amsmath, amssymb, physics}
\usepackage[margin=1in]{geometry}
\usepackage{graphicx}
\usepackage{booktabs}
\usepackage{caption}
\usepackage{float}
\usepackage{amsmath}
\usepackage{enumitem}
\setlength{\parindent}{0pt}
\setlength{\parskip}{6pt}

\newenvironment{solution}
{\par\noindent\textbf{Solution:}\par}
{\vspace{1em}}

\title{Homework Assignment No.3 \\ EE520/EE459 \\ Yuyao Huang}
\date{Due 11:59 PM, March 12, 2026}
\begin{document}
\maketitle

\textbf{Total: 100 points} \\
\textbf{Submit: a single PDF file of all your work}

\section*{Problem 1 (15 points)}
Inspect and convince yourself of the equality of the matrix and dyad representations below.
\[\begin{bmatrix}1 & 0 & 0 & 0 \\ 0 & 1 & 0 & 0 \\ 0 & 0 & a & b \\ 0 & 0 & c & d \end{bmatrix} = \ket{00}\bra{00} + \ket{01}\bra{01} + a\ket{10}\bra{10} + b\ket{10}\bra{11} + c\ket{11}\bra{10} + d\ket{11}\bra{11}\]
\underline{Question:} Write down the dyad form for the matrix below:
\[\begin{bmatrix}1 & 0 & 0 & 0 \\ 0 & a & b & 0 \\ 0 & c & d & 0 \\ 0 & 0 & 0 & 1 \end{bmatrix} =\]

\begin{solution}
Since $\ket{x}\bra{y}$ represents a operator, and its function is:\[\ket{x}\bra{y}\ket{\psi} = \ket{x}\left(\bra{y}\ket{\psi}\right)\]

So,
\begin{itemize}[itemsep=0em]
    \item Row index decides the ket part, and column index decides the bra part.
    \item In the dyad form, the ket part is the row index, and the bra part is the column index.
\end{itemize}

Therefore, the $(i,j)$th element of the matrix corresponds exactly to: \[\ket{i \text{th basis}}\bra{j \text{th basis}}\]

And \begin{align*}
    1 &\leftrightarrow \ket{00} \\ 
    2 &\leftrightarrow \ket{01} \\ 
    3 &\leftrightarrow \ket{10} \\ 
    4 &\leftrightarrow \ket{11}
\end{align*}

Thus \[\begin{bmatrix}(1,1) & (1,2) & (1,3) & (1,4) \\ (2,1) & (2,2) & (2,3) & (2,4) \\ (3,1) & (3,2) & (3,3) & (3,4) \\ (4,1) & (4,2) & (4,3) & (4,4)\end{bmatrix}\]

which corresponds to \[\begin{bmatrix}\ket{00}\bra{00} & \ket{00}\bra{01} & \ket{00}\bra{10} & \ket{00}\bra{11} \\ \ket{01}\bra{00} & \ket{01}\bra{01} & \ket{01}\bra{10} & \ket{01}\bra{11} \\ \ket{10}\bra{00} & \ket{10}\bra{01} & \ket{10}\bra{10} & \ket{10}\bra{11} \\ \ket{11}\bra{00} & \ket{11}\bra{01} & \ket{11}\bra{10} & \ket{11}\bra{11}\end{bmatrix}\]

So now,
\[\begin{bmatrix}1 & 0 & 0 & 0 \\ 0 & a & b & 0 \\ 0 & c & d & 0 \\ 0 & 0 & 0 & 1 \end{bmatrix} = \ket{00}\bra{00} + a\ket{01}\bra{01} + b\ket{01}\bra{10} + c\ket{10}\bra{01} + d\ket{10}\bra{10} + \ket{11}\bra{11}\]

\end{solution}

\section*{Problem 2 (15 points)}
Recall the $+1$ eigenstate $\ket{\Psi_n}$ of $\operatorname{\sigma_n} = \mathbf{n} \cdot \operatorname{\sigma}$ where $\operatorname{\sigma_n}\ket{\Psi_n} = + \ket{\Psi_n}$, where $\mathbf{n}$ is the unit vector with angle coordinates of ($\theta$, $\phi$), which is given by
\[\ket{\psi_n} = \ket{\psi(\theta, \phi)} = \cos{\frac{\theta}{2}}\ket{0} + e^{i\phi}\sin{\frac{\theta}{2}}\ket{1} = \begin{bmatrix}\cos{\frac{\theta}{2}} \\ e^{i\phi}\sin{\frac{\theta}{2}}\end{bmatrix}\]

\underline{Question:} By writing the below two expressions as matrices, show the equality. Here, $P+$ is the projector to the $+1$ eigenstate: \[P+ = \ket{\psi_n}\bra{\psi_n} = \frac{1}{2}\left(\operatorname{I} + \mathbf{n} \cdot \operatorname{\sigma}\right)\]

\begin{solution}

\textbf{Step 1: calculate leftside of the equation}

From the question, we have \[\ket{\psi_n} = \begin{bmatrix}\cos{\frac{\theta}{2}} \\ e^{i\phi}\sin{\frac{\theta}{2}}\end{bmatrix}\]

And the formation of Conjugate Transpose (Hermitian), \[\bra{\psi_n} = \left(\ket{\psi_n}\right)^\dagger = \left(\ket{\psi_n}^*\right)^{\mathrm{T}} = \left(\begin{bmatrix}\cos{\frac{\theta}{2}} \\ e^{-i\phi}\sin{\frac{\theta}{2}}\end{bmatrix}\right)^{\mathrm{T}} = \begin{bmatrix}\cos{\frac{\theta}{2}} & e^{-i\phi}\sin{\frac{\theta}{2}}\end{bmatrix}\]

So, 
\[\begin{aligned}
\ket{\psi_n}\bra{\psi_n} &= \begin{bmatrix}\cos{\frac{\theta}{2}} \\ e^{i\phi}\sin{\frac{\theta}{2}}\end{bmatrix} \begin{bmatrix}\cos{\frac{\theta}{2}} & e^{-i\phi}\sin{\frac{\theta}{2}}\end{bmatrix} \\
&= \begin{bmatrix}\left(\cos{\frac{\theta}{2}}\right)^2 & \cos{\frac{\theta}{2}} \cdot e^{-i\phi}\sin{\frac{\theta}{2}} \\ e^{i\phi}\sin{\frac{\theta}{2}} \cdot \cos{\frac{\theta}{2}} & e^{i\phi}\sin{\frac{\theta}{2}} \cdot e^{-i\phi}\sin{\frac{\theta}{2}}\end{bmatrix} \\
&= \boxed{\begin{bmatrix}\cos^2{\frac{\theta}{2}} & e^{-i\phi}\cos{\frac{\theta}{2}}\sin{\frac{\theta}{2}} \\ e^{i\phi}\cos{\frac{\theta}{2}}\sin{\frac{\theta}{2}} & \sin^2{\frac{\theta}{2}}\end{bmatrix}}
\end{aligned}\]

\textbf{Step 2: calculate rightside of the equation}

$\mathbf{n}$ is the unit vector with angle coordinates of ($\theta$, $\phi$), so we can write $\mathbf{n}$ as \[\mathbf{n} = \left(\sin{\theta}\cos{\phi}, \sin{\theta}\sin{\phi}, \cos{\theta}\right)\]

Pauli matrices are defined as:\[\sigma_x = \begin{bmatrix}0 & 1 \\ 1 & 0\end{bmatrix}, \quad \sigma_y = \begin{bmatrix}0 & -i \\ i & 0\end{bmatrix}, \quad \sigma_z = \begin{bmatrix}1 & 0 \\ 0 & -1\end{bmatrix}\]

Then,
\[\begin{aligned}
\mathbf{n} \cdot \operatorname{\sigma} &= n_x\sigma_x + n_y\sigma_y + n_z\sigma_z \\
&= \left(\sin{\theta}\cos{\phi}\right)\sigma_x + \left(\sin{\theta}\sin{\phi}\right)\sigma_y + \left(\cos{\theta}\right)\sigma_z \\
&= \left(\sin{\theta}\cos{\phi}\right)\begin{bmatrix}0 & 1 \\ 1 & 0\end{bmatrix} + \left(\sin{\theta}\sin{\phi}\right)\begin{bmatrix}0 & -i \\ i & 0\end{bmatrix} + \left(\cos{\theta}\right)\begin{bmatrix}1 & 0 \\ 0 & -1\end{bmatrix} \\
&= \begin{bmatrix}0 & \sin{\theta}\cos{\phi} \\ \sin{\theta}\cos{\phi} & 0\end{bmatrix} + \begin{bmatrix}0 & -i\sin{\theta}\sin{\phi} \\ i\sin{\theta}\sin{\phi} & 0\end{bmatrix} + \begin{bmatrix}\cos{\theta} & 0 \\ 0 & -\cos{\theta}\end{bmatrix} \\
&= \begin{bmatrix}\cos{\theta} & \sin{\theta}\left(\cos{\phi} - i\sin{\phi}\right) \\ \sin{\theta}\left(\cos{\phi} + i\sin{\phi}\right) & -\cos{\theta}\end{bmatrix} \\
&= \begin{bmatrix}\cos{\theta} & \sin{\theta}e^{-i\phi} \\ \sin{\theta}e^{i\phi} & -\cos{\theta}\end{bmatrix}
\end{aligned}\]

$I$ is the identity matrix, so \[\operatorname{I} = \begin{bmatrix}1 & 0 \\ 0 & 1\end{bmatrix}\]

Thus, 
\[\begin{aligned}
\frac{1}{2}\left(\operatorname{I} + \mathbf{n} \cdot \operatorname{\sigma}\right) &= \frac{1}{2} \left( \begin{bmatrix}1 & 0 \\ 0 & 1\end{bmatrix} + \begin{bmatrix}\cos{\theta} & \sin{\theta}e^{-i\phi} \\ \sin{\theta}e^{i\phi} & -\cos{\theta}\end{bmatrix} \right) \\
&= \frac{1}{2} \begin{bmatrix}1 + \cos{\theta} & \sin{\theta}e^{-i\phi} \\ \sin{\theta}e^{i\phi} & 1 - \cos{\theta}\end{bmatrix} \\
&= \begin{bmatrix}\frac{1 + \cos{\theta}}{2} & \frac{\sin{\theta}e^{-i\phi}}{2} \\ \frac{\sin{\theta}e^{i\phi}}{2} & \frac{1 - \cos{\theta}}{2}\end{bmatrix}
\end{aligned}\]

Now we have \[\cos^2{\frac{\theta}{2}} = \frac{1+\cos{\theta}}{2}, \quad \sin^2{\frac{\theta}{2}} = \frac{1-\cos{\theta}}{2}, \quad \sin{\theta} = 2\sin{\frac{\theta}{2}}\cos{\frac{\theta}{2}}\]

So 
\[\frac{\sin{\theta}}{2} = \sin{\frac{\theta}{2}}\cos{\frac{\theta}{2}}\]

\[\boxed{\Rightarrow \frac{1}{2}\left(\operatorname{I} + \mathbf{n} \cdot \operatorname{\sigma}\right) = \begin{bmatrix}\cos^2{\frac{\theta}{2}} & e^{-i\phi}\cos{\frac{\theta}{2}}\sin{\frac{\theta}{2}} \\ e^{i\phi}\cos{\frac{\theta}{2}}\sin{\frac{\theta}{2}} & \sin^2{\frac{\theta}{2}}\end{bmatrix}}\]

\textbf{Step 3: compare the leftside and rightside of the equation}

\[\boxed{\text{We can see that the leftside and rightside of the equation are equal, so the equation is verified.}}\]

\end{solution}

\section*{Problem 3 (25 points)}
Generalization of Bloch sphere picture (Sec. 1.2) to the \textbf{\textit{mixed}} state for single qubits.
\begin{enumerate}
    \item \underline{Question:} Show that for a \textbf{pure} state, the density matrix is $\rho = \frac{1}{2}(I + \mathbf{n} \cdot \boldsymbol{\sigma})$ where $\mathbf{n}$ is a unit vector. [Sec.1.2 says that $\mathbf{r}$ for $\rho = \ket{\psi(\theta,\phi)}\bra{\psi(\theta,\phi)}$ is a unit vector with angular coordinates $(\theta, \phi)$ on the Bloch sphere.] \textit{Hint}: A pure state may be considered an ensemble $\{p=1, \ket{\psi_n}\}$ of a single state $\ket{\psi_n}$ and probability $p=1$.
    \item \underline{Question:} Show that the \textbf{\textit{density matrix}} for a mixed state given by an ensemble of pure states $\{p_j, \ket{\psi_j}\}$ where $\sum_j p_j = 1$ may be written as $\rho = \dfrac{I + \mathbf{r} \cdot \boldsymbol{\sigma}}{2}$ where $\mathbf{r}$ is a 3-dimensional vector such that $\|\mathbf{r}\| \leq 1$. This $\mathbf{r}$ is known as the \textbf{\textit{Bloch vector}} for $\rho$. \textit{Hint}: $\rho = \sum_j p_j \ket{\psi_j}\bra{\psi_j} = \sum_j p_j \frac{1}{2}(I + \mathbf{n}_j \cdot \boldsymbol{\sigma}) = \frac{1}{2}\left[I\sum_j p_j + \sum_j p_j \mathbf{n}_j \cdot \boldsymbol{\sigma}\right]$ $\Rightarrow$ find the expression for the Bloch vector $\mathbf{r}$ and argue why the size of $\mathbf{r}$ will be less than or equal to 1. The equality is when all pure state Bloch vectors $\mathbf{n}_j$ are aligned on a straight line in the same direction to the same point $\mathbf{n}$ on the Bloch sphere. In this case, $\sum_j p_j \mathbf{n}_j = \mathbf{n}$ for simply one pure state.
    \item \underline{Question:} What is the \textbf{\textit{Bloch vector}} $\mathbf{r}$ for $\rho = I/2$?
    \item \underline{Question:} Show that the state is \textbf{pure} if and only if $\|\mathbf{r}\| = 1$. [\textit{Hint}: Apply $\text{Tr}(\rho^2) = 1$, $(\mathbf{r} \cdot \boldsymbol{\sigma})^2 = r^2 I$, anti-commuting $\boldsymbol{\sigma}$ and $\boldsymbol{\sigma}^2 = I$ to show $r^2 = 1$.]
\end{enumerate}

\begin{solution}

\textbf{Solution 1}

For a pure state, the density matrix is defined as\[\rho = \ket{\psi_n}\bra{\psi_n}.\]
    
A pure state may be viewed as an ensemble consisting of a single state \(\{p=1, \ket{\psi_n}\}\).
    
Hence, \[\rho = 1\cdot \ket{\psi_n}\bra{\psi_n}\]
    
From Problem 2, we have already shown that \[\ket{\psi_n}\bra{\psi_n}=\frac12 \left(I+\mathbf n\cdot\boldsymbol{\sigma}\right)\]

where \(\mathbf n=(\sin\theta\cos\phi,\sin\theta\sin\phi,\cos\theta)\) is a unit vector corresponding to the point \((\theta,\phi)\) on the Bloch sphere.

Therefore, the density matrix for a pure qubit state can be written as \[\rho=\frac12\left(I+\mathbf n\cdot\boldsymbol{\sigma}\right)\]
where \(\boxed{\|\mathbf n\|=1}\).

\textbf{Solution 2}

For a mixed state described by an ensemble of pure states $\{p_j, \ket{\psi_j}\}$ with $\sum_j p_j=1$, the density matrix is \[\rho=\sum_j p_j \ket{\psi_j}\bra{\psi_j}.\]

From part (1), each pure-state density matrix can be written as \[\ket{\psi_j}\bra{\psi_j} = \frac12\left(I+\mathbf n_j\cdot\boldsymbol{\sigma}\right),\]

where $\mathbf n_j$ is a unit Bloch vector, so $\|\mathbf n_j\|=1$. Substituting into $\rho$, we obtain \[\rho = \sum_j p_j \frac12\left(I+\mathbf n_j\cdot\boldsymbol{\sigma}\right) = \frac12\left[I\sum_j p_j+\sum_j p_j\,\mathbf n_j\cdot\boldsymbol{\sigma}\right].\]

Since $\sum_j p_j=1$, this becomes \[\rho = \frac12\left[I+\sum_j p_j\,\mathbf n_j\cdot\boldsymbol{\sigma}\right].\]

Define the Bloch vector $\mathbf r=\sum_j p_j \mathbf n_j$, then \[\boxed{\rho=\frac12\left(I+\mathbf r\cdot\boldsymbol{\sigma}\right)}.\]

Now we show $\|\mathbf r\|\le 1$. By the triangle inequality, \[\|\mathbf r\| = \left\|\sum_j p_j \mathbf n_j\right\| \le \sum_j p_j \|\mathbf n_j\| = \sum_j p_j = 1.\]

Therefore $\boxed{\|\mathbf r\|\le 1}$, and the Bloch vector of any single-qubit mixed state satisfies $\|\mathbf r\|\le 1$.

\textbf{Solution 3}

From part (2), we have \[\rho=\dfrac12\left(\operatorname{I}+\mathbf{r}\cdot\boldsymbol{\sigma}\right)\]

Comparing with \(\rho=\frac{\operatorname{I}}{2}\), we have \[\frac{\operatorname{I}}{2} = \dfrac12\left(\operatorname{I}+\mathbf{r}\cdot\boldsymbol{\sigma}\right) = \frac{\operatorname{I}}{2} + \frac12\left(\mathbf{r}_x\sigma_x+\mathbf{r}_y\sigma_y+\mathbf{r}_z\sigma_z\right)\]

gives \[r_x\sigma_x+r_y\sigma_y+r_z\sigma_z=0\]

Multiplying by \(\sigma_k\) and taking the trace, and using the orthogonality \[\operatorname{Tr}(\sigma_i\sigma_j)=2\delta_{ij}\]

yields \(2r_k=0\) for each \(k=x,y,z\).

Therefore, \[\boxed{\mathbf{r}=(0,0,0)}\]

i.e., \(\rho=I/2\) corresponds to the center of the Bloch sphere.

\textbf{Solution 4}

For a single qubit, \[\rho=\frac12\left(I+\mathbf r\cdot\boldsymbol{\sigma}\right).\]

We compute \[\rho^2 = \frac14\left(I+\mathbf r\cdot\boldsymbol{\sigma}\right)^2 = \frac14\left(I+2\,\mathbf r\cdot\boldsymbol{\sigma}+(\mathbf r\cdot\boldsymbol{\sigma})^2\right)\]

Using the Pauli matrix relations $\sigma_x^2=\sigma_y^2=\sigma_z^2=I$ and $\sigma_i\sigma_j+\sigma_j\sigma_i=0$ for $i\ne j$, we get \[(\mathbf r\cdot\boldsymbol{\sigma})^2 = (r_x^2+r_y^2+r_z^2)I = \|\mathbf r\|^2 I\]

Therefore, \[\rho^2 = \frac14\left[(1+\|\mathbf r\|^2)I+2\,\mathbf r\cdot\boldsymbol{\sigma}\right]\]

Taking the trace, and using $\operatorname{Tr}(I)=2$, $\operatorname{Tr}(\sigma_x)=\operatorname{Tr}(\sigma_y)=\operatorname{Tr}(\sigma_z)=0$, it follows that \[\operatorname{Tr}(\rho^2) = \frac12(1+\|\mathbf r\|^2)\]

A state is pure if and only if $\operatorname{Tr}(\rho^2)=1$. Thus, \[1=\frac12(1+\|\mathbf r\|^2) \quad\Longrightarrow\quad \|\mathbf r\|^2=1 \quad\Longrightarrow\quad \|\mathbf r\|=1\]

So, if $\rho$ is pure, then $\|\mathbf r\|=1$. Conversely, if $\|\mathbf r\|=1$, then $(\mathbf r\cdot\boldsymbol{\sigma})^2=I$, and therefore \[\rho^2 = \frac14\left(I+2\,\mathbf r\cdot\boldsymbol{\sigma}+I\right) = \frac12\left(I+\mathbf r\cdot\boldsymbol{\sigma}\right) = \rho\]

Hence $\rho$ is pure. Therefore, \[\boxed{\rho \text{ is pure if and only if } \|\mathbf r\|=1}\]

\end{solution}

\section*{Problem 4 (20 points)}
Find the Schmidt decomposition of the states below. [Important: Make sure that your Schmidt basis is orthonormal. If not, then your answer is wrong.]
\[\frac{1}{\sqrt{2}}(\ket{00} + \ket{11}), \quad \frac{2\ket{00} + \ket{01} + \ket{10} + \ket{11}}{\sqrt{10}}\]

\begin{solution}

\textbf{(1) State 1}

The state $\ket{\psi_1}=\dfrac{1}{\sqrt{2}}(\ket{00}+\ket{11})$ is already in Schmidt form:
\[
\ket{\psi_1}=\frac{1}{\sqrt{2}}\ket{0}\ket{0}+\frac{1}{\sqrt{2}}\ket{1}\ket{1}.
\]

The Schmidt coefficients are $\lambda_1=\lambda_2=\dfrac{1}{\sqrt{2}}$, with Schmidt bases $\ket{u_1}=\ket{v_1}=\ket{0}$ and $\ket{u_2}=\ket{v_2}=\ket{1}$. These are orthonormal since $\langle0|0\rangle=\langle1|1\rangle=1$ and $\langle0|1\rangle=0$. Therefore,
\[
\boxed{\ket{\psi_1}=\frac{1}{\sqrt{2}}\ket{0}\ket{0}+\frac{1}{\sqrt{2}}\ket{1}\ket{1}},
\]
which is the standard Schmidt decomposition of the Bell state.

\textbf{(2) State 2}

For a bipartite state $\ket{\psi}=\sum_{i,j}A_{ij}\ket{i}\ket{j}$, the Schmidt decomposition is obtained via SVD of the coefficient matrix $A$. From $\ket{\psi_2}=\dfrac{2\ket{00}+\ket{01}+\ket{10}+\ket{11}}{\sqrt{10}}$, we read off
\[A=\frac{1}{\sqrt{10}}\begin{bmatrix}2&1\\1&1\end{bmatrix}\]

\textbf{Schmidt coefficients.} The squared Schmidt coefficients are eigenvalues of $AA^\dagger$:
\[AA^\dagger = \frac{1}{\sqrt{10}}\begin{bmatrix}2&1\\1&1\end{bmatrix} \frac{1}{\sqrt{10}}\begin{bmatrix}2&1\\1&1\end{bmatrix}^{\mathrm{T}} = \frac{1}{10}\begin{bmatrix}5&3\\3&2\end{bmatrix}\]

The eigenvalues of $\begin{bmatrix}5&3\\3&2\end{bmatrix}$ satisfy $\det(M-\mu I)=(5-\mu)(2-\mu)-9=\mu^2-7\mu+1=0$, giving $\mu_{1,2}=\dfrac{7\pm3\sqrt5}{2}$. 
Hence
\[\lambda_{1,2}^2=\frac{7\pm3\sqrt5}{20}, \qquad \boxed{\lambda_1=\sqrt{\frac{7+3\sqrt5}{20}},\quad\lambda_2=\sqrt{\frac{7-3\sqrt5}{20}}}.\]

\textbf{Schmidt basis for the first subsystem.} Letting $\varphi=\dfrac{1+\sqrt5}{2}$, the eigenvectors of $M$ are found by solving $(M-\mu_i I)\mathbf{x}=0$. For $\mu_1$: $y=\dfrac{\sqrt5-1}{2}x=\dfrac{1}{\varphi}x$, giving the unnormalized vector $(\varphi,1)^{\mathrm{T}}$. For $\mu_2$: $y=-\varphi x$, giving $(1,-\varphi)^{\mathrm{T}}$. Normalizing and verifying orthogonality $\langle u_1|u_2\rangle\propto\varphi\cdot1+1\cdot(-\varphi)=0$:
\[
\ket{u_1}=\frac{\varphi\ket0+\ket1}{\sqrt{\varphi^2+1}},
\qquad
\ket{u_2}=\frac{\ket0-\varphi\ket1}{\sqrt{\varphi^2+1}}.
\]

\textbf{Schmidt basis for the second subsystem.} Since $A$ is real symmetric, the left and right singular vectors coincide, so we take $\ket{v_i}=\ket{u_i}$.

\textbf{Schmidt decomposition:}
\[
\boxed{\ket{\psi_2}=\sqrt{\frac{7+3\sqrt5}{20}}\,\ket{u_1}\ket{v_1}+\sqrt{\frac{7-3\sqrt5}{20}}\,\ket{u_2}\ket{v_2}},
\]
where $\ket{u_1}=\ket{v_1}=\dfrac{\varphi\ket0+\ket1}{\sqrt{\varphi^2+1}}$, $\ket{u_2}=\ket{v_2}=\dfrac{\ket0-\varphi\ket1}{\sqrt{\varphi^2+1}}$, and $\varphi=\dfrac{1+\sqrt5}{2}$.

\end{solution}

\section*{Problem 5 (25 points)}
Define \(k = k_{n-1}2^{n-1} + k_{n-2}2^{n-2} + \cdots + k_0 2^0\), \(\ket{k} = \ket{k_{n-1}}\ket{k_{n-2}}\cdots\ket{k_0}\), where \(k\) is an integer ranging from 0 to \(2^n-1\) and \(k_j \in \{0,1\}\). Assume \(\theta\) is a constant parameter. Show that

\[\sum_{k=0}^{2^n-1} e^{i\theta k}\ket{k} = \left(\ket{0} + e^{i\theta 2^{n-1}}\ket{1}\right) \otimes \left(\ket{0} + e^{i\theta 2^{n-2}}\ket{1}\right) \otimes \cdots \otimes \left(\ket{0} + e^{i\theta 2^0}\ket{1}\right)\]

\underline{Note:} The above result shows that \textbf{LHS = sum of \(2^n\) terms \(\boldsymbol{\rightarrow}\) RHS = tensor product of \(\textit{n}\) terms}. This leads to the exponential speed up of the Quantum Fourier Transform (QFT) over the classical Fourier transform.

\underline{Hint:} Note that the LHS of the integer summation is broken into the RHS of tensor products of binary summations. For the RHS below, factor it into the different \(k_j\) terms and explicitly do the sum for \(k_j = 0 \text{ and } 1\)

\[\sum_{k=0}^{2^n-1} e^{i\theta k}\ket{k} = \sum_{k_{n-1},k_{n-2},\cdots,k_0=0}^1 e^{i\theta \left(k_{n-1}2^{n-1} + k_{n-2}2^{n-2} + \cdots + k_0 2^0\right)}\ket{k_{n-1}k_{n-2}\cdots k_0}\]

\begin{solution}

\textbf{Step 1: Express \(k\) in binary form.}

By definition, $k=k_{n-1}2^{n-1}+\cdots+k_0 2^0$ with $k_j\in\{0,1\}$, and $\ket{k}=\ket{k_{n-1}}\otimes\cdots\otimes\ket{k_0}$. Since each $k_j$ independently takes values in $\{0,1\}$, summing over $k$ from $0$ to $2^n-1$ is equivalent to summing over all binary digits independently:
\[
\sum_{k=0}^{2^n-1} e^{i\theta k}\ket{k}
=
\sum_{k_{n-1},\dots,k_0=0}^1
e^{i\theta\left(k_{n-1}2^{n-1}+\cdots+k_0 2^0\right)}
\left(\ket{k_{n-1}}\otimes\cdots\otimes\ket{k_0}\right).
\]

\textbf{Step 2: Factor the exponential.}

Using $e^{A+B}=e^Ae^B$, the exponential factors as $e^{i\theta k_{n-1}2^{n-1}}\cdots e^{i\theta k_0 2^0}$. Distributing each scalar into the corresponding tensor factor gives
\[
\sum_{k=0}^{2^n-1} e^{i\theta k}\ket{k}
=
\sum_{k_{n-1},\dots,k_0=0}^1
\left(e^{i\theta k_{n-1}2^{n-1}}\ket{k_{n-1}}\right)
\otimes\cdots\otimes
\left(e^{i\theta k_0 2^0}\ket{k_0}\right).
\]

\textbf{Step 3: Separate into independent single-bit sums.}

Since each term in the multiple sum is a tensor product where the $j$-th factor $e^{i\theta k_j 2^j}\ket{k_j}$ depends only on $k_j$, and the variables $k_0,\dots,k_{n-1}$ are independent, we can apply the distributivity of the tensor product over addition, i.e.\ $(A\otimes B)+(A\otimes C)=A\otimes(B+C)$, repeatedly to factor the multiple sum into a tensor product of single-variable sums:
\[
\sum_{k=0}^{2^n-1} e^{i\theta k}\ket{k}
=
\left(\sum_{k_{n-1}=0}^1 e^{i\theta k_{n-1}2^{n-1}}\ket{k_{n-1}}\right)
\otimes\cdots\otimes
\left(\sum_{k_0=0}^1 e^{i\theta k_0 2^0}\ket{k_0}\right).
\]

\textbf{Step 4: Evaluate each single-bit sum.}

For each factor, since $k_j\in\{0,1\}$ and $e^0=1$:
\[
\sum_{k_j=0}^1 e^{i\theta k_j 2^j}\ket{k_j}
= e^0\ket{0}+e^{i\theta 2^j}\ket{1}
= \ket{0}+e^{i\theta 2^j}\ket{1}.
\]

Applying this to every factor yields the desired result:
\[
\boxed{
\sum_{k=0}^{2^n-1} e^{i\theta k}\ket{k}
=
\left(\ket{0}+e^{i\theta 2^{n-1}}\ket{1}\right)\otimes
\left(\ket{0}+e^{i\theta 2^{n-2}}\ket{1}\right)\otimes\cdots\otimes
\left(\ket{0}+e^{i\theta 2^0}\ket{1}\right)
}.
\]

This shows that a sum of $2^n$ terms on the LHS equals a tensor product of only $n$ factors on the RHS, which is the key structural identity underlying the exponential speedup of the Quantum Fourier Transform.

\end{solution}

\end{document}